\documentclass[conference]{IEEEtran}
\IEEEoverridecommandlockouts
\usepackage{cite}
\usepackage{amsmath,amssymb,amsfonts}
\usepackage{graphicx}
\usepackage{textcomp}
\usepackage{xcolor}
\usepackage{url}
\usepackage{hyperref}
\usepackage{booktabs} % Para mejores tablas
% TeX Live ``basic'' a veces no incluye Courier (pcr); usar cmtt para \texttt
\renewcommand{\ttdefault}{cmtt}
\def\BibTeX{{\rm B\kern-.05em{\sc i\kern-.025em b}\kern-.08em
    T\kern-.1667em\lower.7ex\hbox{E}\kern-.125emX}}
\begin{document}

\title{Proyecto I: Reconocimiento de comandos de voz y
despliegue móvil con ONNX Runtime Mobile}

\author{\IEEEauthorblockN{Fabricio Mena Mejía}
\IEEEauthorblockA{Instituto Tecnológico de Costa Rica\\
Email: fabriciomena11@estudiantec.cr}
\and
\IEEEauthorblockN{Anthony Fuentes Calvo}
\IEEEauthorblockA{Instituto Tecnológico de Costa Rica\\
Email: anthonyfuentescalvo@estudiantec.cr}
\and
\IEEEauthorblockN{Fernando Sánchez Hidalgo}
\IEEEauthorblockA{Instituto Tecnológico de Costa Rica\\
Email: Fer.sanchez@estudiantec.cr}}

\maketitle

\begin{abstract}

\end{abstract}

\section{Marco teórico}

\section{Diseño de arquitectura de IA}

\subsection{Control de Aleatoriedad y Reproducibilidad}
Para garantizar la reproducibilidad de los experimentos, se implementó un estricto control de aleatoriedad. Se definió una semilla global (\textit{seed} $= 42$) utilizada para inicializar los generadores de números pseudoaleatorios de las librerías \texttt{random}, \texttt{numpy} y \texttt{torch}. Adicionalmente, se configuró el comportamiento determinista de cuDNN desactivando el \textit{benchmark} dinámico donde fuera necesario.

\subsection{Modelo A: Datos Crudos (LeNet-5 Adaptado)}
El Modelo A está basado en la arquitectura clásica de LeNet-5, pero fue adaptado específicamente para el procesamiento de mel-espectrogramas. 

\subsubsection{Justificación de modificaciones}
Se realizaron los siguientes cambios sobre la arquitectura original de LeCun:
\begin{itemize}
    \item \textbf{Tamaño de entrada:} Se modificó para recibir tensores de $1 \times 128 \times 128$ correspondientes a las imágenes generadas del audio, en lugar de los $32 \times 32$ originales.
    \item \textbf{Filtros convolucionales:} Se aumentó el número de canales a 32, 64 y 128 en las tres capas convolucionales consecutivas para capturar patrones espectrales más ricos y complejos.
    \item \textbf{Función de Activación:} Se preserva el uso de la función de activación clásica \texttt{Tanh} estipulada en la arquitectura original de LeNet-5. Esto se hace en estricto cumplimiento con los requerimientos del proyecto que limitan el uso de funciones más modernas como \texttt{ReLU}, garantizando así que el modelo mantenga la esencia matemática clásica de la red mientras se adapta al dominio de audio.
    \item \textbf{Padding:} En las capas convolucionales iniciales se utiliza un \texttt{padding=0}, consistentemente con la versión clásica de LeNet. Esta elección genera una reducción espacial después de cada convolución antes del \texttt{AvgPool2d}, lo que simplifica el diseño y mantiene la estructura de extracción de características sin introducir ceros adicionales en los bordes.
    \item \textbf{Clasificador 100\% Convolucional:} Se reemplazó la sección tradicional de capas totalmente conectadas (\textit{Fully Connected}) por una operación de \textit{Global Average Pooling} seguida de convoluciones $1 \times 1$. Esta modificación reduce significativamente el número de parámetros y mitiga el sobreajuste (\textit{overfitting}), preparando el modelo para su futuro despliegue en dispositivos móviles donde los parámetros en capas densas consumen demasiada memoria RAM.
\end{itemize}

\subsubsection{Definición de la Arquitectura}
La arquitectura final se resume en el diagrama presentado en la Tabla \ref{tab:arquitectura}. Consta de tres bloques convolucionales iniciales con \textit{Average Pooling}, seguidos de una capa de \textit{Global Average Pooling} y un bloque clasificador compuesto exclusivamente por convoluciones de $1 \times 1$.

\begin{table}[htbp]
\caption{Diagrama de Capas - LeNet-5 Adaptado}
\label{tab:arquitectura}
\centering
\begin{tabular}{ll}
\toprule
\textbf{Capa / Operación} & \textbf{Dimensión de Salida} \\
\midrule
Input (Mel-Spectrogram) & $(B, 1, 128, 128)$ \\
Conv1 ($32 \times 5 \times 5$) + Tanh + AvgPool(2) & $(B, 32, 62, 62)$ \\
Conv2 ($64 \times 5 \times 5$) + Tanh + AvgPool(2) & $(B, 64, 29, 29)$ \\
Conv3 ($128 \times 5 \times 5$) + Tanh + AvgPool(2) & $(B, 128, 12, 12)$ \\
Global Average Pooling & $(B, 128, 1, 1)$ \\
Conv4 ($128 \to 512, 1 \times 1$) + Tanh & $(B, 512, 1, 1)$ \\
Conv5 ($512 \to 128, 1 \times 1$) + Tanh & $(B, 128, 1, 1)$ \\
Conv6 ($128 \to 10, 1 \times 1$) (Logits) & $(B, 10, 1, 1)$ \\
Flatten Final & $(B, 10)$ \\
\bottomrule
\end{tabular}
\end{table}

\begin{figure}[htbp]
\centerline{\includegraphics[width=\linewidth]{images/Modelo_LeNet5_adaptado.png}}
\caption{Diagrama visual de la arquitectura del Modelo A.}
\label{fig:model_a}
\end{figure}

\subsubsection{Función de Pérdida, Optimizador y Entrenamiento}
Se utilizó \textit{Cross-Entropy Loss} como función de pérdida, la cual es ideal para clasificación multiclase. El optimizador seleccionado fue \texttt{Adam} con una tasa de aprendizaje de $1 \times 10^{-3}$ y \textit{weight decay} de $1 \times 10^{-4}$ para mitigar el sobreajuste. Además, se empleó un \textit{scheduler} (\texttt{ReduceLROnPlateau}) que reduce la tasa de aprendizaje a la mitad si la precisión de validación no mejora en 3 épocas. Las rutinas de entrenamiento se ejecutaron usando precisión mixta nativa (\texttt{torch.amp}).


\section{Análisis y comparación de resultados}

\subsection{Modelo A: Datos crudos}

El Modelo A actual —con la arquitectura LeNet-5 adaptada a mel-espectrogramas, activaciones \texttt{Tanh} y un clasificador 100\% convolucional— alcanzó los siguientes resultados en la evaluación final del conjunto de prueba:

\begin{itemize}
    \item \textbf{Loss (test):} 1.0522
    \item \textbf{Accuracy (test):} 0.6388
    \item \textbf{F1 Macro (test):} 0.6344
    \item \textbf{Mejor Accuracy de validación:} 0.6300
\end{itemize}

Estos valores reflejan el comportamiento del modelo actual con datos crudos, el desglose de resultados para cada comando de voz se muestra en la tabla \ref{tab:metricas_a}

\begin{table}[htbp]
\caption{Reporte de Clasificación - Modelo A (datos crudos)}
\label{tab:metricas_a}
\centering
\begin{tabular}{lcccc}
\toprule
\textbf{Comando} & \textbf{Precision} & \textbf{Recall} & \textbf{F1-Score} & \textbf{Support} \\
\midrule
yes   & 0.87 & 0.57 & 0.69 & 630 \\
no    & 0.51 & 0.33 & 0.40 & 569 \\
up    & 0.55 & 0.82 & 0.66 & 573 \\
down  & 0.54 & 0.77 & 0.64 & 586 \\
left  & 0.61 & 0.70 & 0.65 & 555 \\
right & 0.83 & 0.72 & 0.77 & 576 \\
on    & 0.68 & 0.67 & 0.67 & 574 \\
off   & 0.69 & 0.49 & 0.57 & 565 \\
stop  & 0.65 & 0.73 & 0.69 & 558 \\
go    & 0.59 & 0.60 & 0.59 & 597 \\
\midrule
\textbf{Accuracy} & \multicolumn{3}{c}{\textbf{0.64}} & \textbf{5783} \\
\textbf{Macro Avg}& 0.65 & 0.64 & 0.63 & 5783 \\
\bottomrule
\end{tabular}
\end{table}

\begin{figure}[htbp]
\centerline{\includegraphics[width=\linewidth]{images/Grafica_Loss_Modelo_A_raw.png}}
\caption{Pérdida de entrenamiento y validación para el Modelo A.}
\label{fig:loss_a}
\end{figure}

En la gráfica de \textbf{Loss}, la curva de entrenamiento desciende de forma consistente desde aproximadamente 2.19 hasta 1.22, mientras que la pérdida de validación se reduce desde 2.05 hasta 1.24. Este comportamiento indica que el modelo aprende patrones útiles sin incurrir en un sobreajuste abrupto durante la fase final del entrenamiento.

\begin{figure}[htbp]
\centerline{\includegraphics[width=\linewidth]{images/Grafica_Accuracy_Modelo_A_raw.png}}
\caption{Precisión de entrenamiento y validación para el Modelo A.}
\label{fig:accuracy_a}
\end{figure}

La gráfica de \textbf{Accuracy} muestra una tendencia de aprendizaje gradual: la precisión de entrenamiento crece de alrededor de 16\% a 58\%, mientras que la precisión de validación alcanza un pico de 63\% en las primeras épocas y luego se estabiliza cerca de 59\%. Esto sugiere que el modelo captura información discriminativa con el conjunto actual, pero su capacidad de generalización permanece limitada por la complejidad del problema y por el uso de datos crudos sin augmentation.

\begin{figure}[htbp]
\centerline{\includegraphics[width=\linewidth]{images/Grafica_F1_Modelo_A_raw.png}}
\caption{Evolución del F1-Score macro en validación para el Modelo A.}
\label{fig:f1_a}
\end{figure}

La gráfica de \textbf{F1-Score macro} evidencia que el desempeño del modelo mejora rápidamente en las primeras épocas y luego se estabiliza, reflejando un balance moderado entre precisión y recall. El valor final de F1 macro en prueba es 0.6344, lo que confirma que aún existen diferencias sustanciales entre clases y que algunos comandos son más difíciles de distinguir.

\begin{figure}[htbp]
\centerline{\includegraphics[width=\linewidth]{images/Matriz_Confusion_Modelo_A_raw.png}}
\caption{Matriz de confusión del Modelo A (datos crudos) sobre los datos de prueba.}
\label{fig:matriz_a}
\end{figure}

La matriz de confusión muestra que los comandos \texttt{yes} y \texttt{no} presentan confusiones relevantes, al igual que pares como \texttt{up}/\texttt{down} y \texttt{go}/\texttt{stop}. Esto es consistente con el desempeño de la arquitectura actual: las activaciones \texttt{Tanh} y el clasificador 100\% convolucional mantienen un modelo compacto, pero la representación aprendida todavía requiere mayor robustez frente a variaciones acústicas.

El rediseño del clasificador con convoluciones $1 \times 1$ y Global Average Pooling reduce parámetros y favorece el despliegue móvil, aunque la versión actual del modelo demuestra que se necesita complementar esta eficiencia con técnicas adicionales como \textit{data augmentation}, normalización avanzada o ajustes finos de hiperparámetros para aumentar la precisión hacia el 70\% o superior.

\section{Referencias}
\begin{thebibliography}{00}

\end{thebibliography}

\end{document}